\chapter{Opening a file}\label{ch:open}
\begin{goals}
\textbf{In this chapter} you open files, learn how lamina decides what a file is, how it reads delimited
text in all its dialects, and how to correct it when it decides wrong.
\end{goals}

\section{Four ways in}\label{sec:ways-in}

\begin{description}
\item[\menu{File>Open…}] or \key{Ctrl}+\key{O}: the system's file picker. In Chrome and Edge the picker
hands lamina a handle it can keep, so the file can be reopened from the recents later.
\item[Drag and drop] a file anywhere onto the window. Each file opens in a tab of its own, so dropping
several opens them all (\cref{ch:views}).
\item[A recent file] from the \menu{File} menu or the empty screen (\cref{sec:recents}).
\item[The operating system] where lamina is installed as an app, opening a \file{.lamina} or \file{.lam}
lens file with it (\cref{ch:lens}). lamina registers only its own lens extensions, not \file{.csv} or
\file{.dm}: it won't take those over from Excel or Datamine.
\end{description}

\section{How lamina decides what a file is}\label{sec:detection}

lamina doesn't go by the file's extension, which is often wrong or missing. It looks at the bytes, in
this order, and takes the first answer that fits:

\begin{enumerate}
\item \textbf{A lens}: a small JSON file carrying lamina's marker. It is applied to the open table, not
opened as data (\cref{ch:lens}).
\item \textbf{A workbook}: a zip whose name ends \file{.xlsx} or \file{.xlsm} (\cref{sec:xlsx}).
\item \textbf{An archive}: zip, tar, gzip, Zstandard, xz or bzip2, recognised by the signature bytes at
its start (\cref{sec:archives}).
\item \textbf{Apache Parquet}: the four bytes \code{PAR1} at both ends of the file
(\cref{sec:parquet}).
\item \textbf{Datamine \file{.dm}}: a recognisable first page of field definitions (\cref{sec:dm}).
\item \textbf{A file opened before}: if its index is in the cache (same name, size and modification
time), lamina reuses it with the interpretation it found then, and the table appears without a scan
(\cref{sec:index-cache}).
\item \textbf{Text}: anything else is examined as text, from its first 64~KB. A table if a delimiter
splits its lines consistently; plain text otherwise; binary if it holds the zero bytes and control
characters text doesn't.
\end{enumerate}

The kind badge at the top right names the result.

\section{Delimited text}\label{sec:delimited}

\subsection{The delimiter}

lamina tries comma, tab, semicolon, pipe and runs of whitespace, and picks the one that splits the most
lines into the same number of columns. A delimiter that gives a consistent count wins; a file whose name
ends \file{.csv}, \file{.tsv} or \file{.tab} is read as a table even when its rows are ragged, using the
best delimiter found. A file with no consistent delimiter and no such name opens as plain text
(\cref{sec:text}).

Whitespace is the delimiter of scientific dumps: columns of numbers separated by spaces, aligned or not.
lamina considers it only when at least half the tokens are numbers, so that ordinary prose, which also
splits on spaces, isn't taken for a table. A real delimiter on the same data wins over whitespace.

\subsection{The header}

A header row is assumed when every field of the first row is text and some later row has a number in it.
A table that is all text, such as a list of codes and descriptions, has no number to go by; in a file
named \file{.csv}, \file{.tsv} or \file{.tab}, its first row is taken as the header when it is complete
(no empty field), its names are all different, and none of them appears again in its column further down.
Otherwise the names are read as the first data row and the columns are called \col{col 1}, \col{col 2},
and so on: set \ui{Header} to \ui{yes} in the interpretation popover (\cref{sec:interpretation}).

A column with no name in the header is called \col{col} and its number. The number of columns is the
length most rows have, so one stray row with an extra field doesn't add a column (data quality flags
such rows, \cref{ch:dq}).

\subsection{Quoting}

Fields may be quoted with double quotes, as Excel and most tools write them. Inside quotes a delimiter is
part of the field, a line break is part of the field (the record goes on to the next line), and a doubled
quote is a quote: \code{"Smith, J."} is one field, and \code{"said ""hi"""} reads \code{said "hi"}.

\subsection{A comment preamble}

Many geological exports start with a block of metadata lines before the header, each beginning with
\code{\#}. When the first line starts with \code{\#}, lamina skips every line at the start that does, and
reads the header from the first line that doesn't. The status line says so:
\code{2 \#-comment lines skipped}. A preamble marked with another character can be named in the
popover's \ui{Comment} box, and a preamble with no marker at all is skipped with \ui{Skip rows}.

\subsection{GSLIB and Geo-EAS}\label{sec:geoeas}

The GSLIB format \parencite{deutsch1998} names its columns in a preamble rather than a header: a title
line, a line with the number of columns $n$, then $n$ lines each holding one column name, then the data,
usually separated by whitespace. lamina recognises the shape, takes the names from the preamble, and
starts the data after it.

\begin{codeblock}
sample data
3
X
Y
GRADE
1.0 2.0 0.5
2.0 3.0 1.25
\end{codeblock}

\subsection{The decimal comma}\label{sec:decimal-comma}

In Brazil and much of Europe a number is written \code{45,2}, and the column delimiter is then usually a
semicolon. lamina recognises it: when the delimiter isn't a comma, and the first lines hold at least
three numbers written with a comma (\code{45,2}, \code{1.234,56}) and over four times as many as are
written with a point, it reads the file with a decimal comma. \code{45,2} then reads as 45.2 and
\code{1.234,56} as 1234.56 (the point is taken as a thousands separator), and filters and statistics treat
those columns as numbers.

When the guess is wrong, or a file has too few such numbers near its top to tell, set \ui{Decimal} to
\ui{point (1.5)} or \ui{comma (1,5)} in the popover.

\subsection{Encodings}\label{sec:encodings}

lamina reads text as UTF-8 unless told otherwise. A file written by older Windows software is often in
Windows-1252 or Latin-1, where accented letters are single bytes that aren't valid UTF-8. Read as UTF-8,
\code{maçica} shows as \code{ma�ica}.

lamina notices: when bytes in the file's first 64~KB don't decode as UTF-8, a \ui{⚠ encoding?} hint
appears in the status line. Click it, or set \ui{Encoding} in the popover, to \ui{Western
(Windows-1252)}, \ui{Latin-1 (ISO-8859-1)}, \ui{UTF-16 LE} or \ui{UTF-16 BE}. The hint goes once you have
chosen an encoding.

\section{Plain text and binary}\label{sec:text}

A text file that isn't a table opens as one column, \col{line}, with a row per line. You can scroll it,
jump to a line, filter it (\code{line contains "ERROR"}) and export what the filter keeps: a log or a
report of any size is readable this way.

A file that holds bytes text never does opens as a note in its tab, \ui{not a table}, with its size.
\menu{View>Raw — the bytes in hex} shows its bytes (\cref{ch:raw}).

\section{The interpretation popover}\label{sec:interpretation}

Click the kind badge, or open \menu{Data>Interpretation (delimiter / header / skip)…}. The popover lists
what lamina detected, under \ui{Detected}, shows the file's first lines as lamina reads them, and gives
the settings that override what it detected.

\shot[0.7\linewidth]{interpretation}{The interpretation popover on a small CSV: what lamina detected, the
first lines as they're read, and the settings.}

\subsection{The first lines, as they're read}\label{sec:first-lines}

The pane headed \ui{the first lines, as they're read} shows the start of the file's own text, line by
line, marked with how lamina is reading it:

\begin{itemize}
\item each delimiter that splits fields is marked; one inside quotes, which doesn't, isn't;
\item the preamble lines skipped, and comment lines, are struck through;
\item the header line is labelled as the header;
\item beside each record, its number of fields, in red where it differs from the table's;
\item a quoted field that runs on to the next line is shown continuing, as one record;
\item the line endings, and invisible characters (tabs, a byte-order mark, a no-break space), as visible
marks, as in the raw view (\cref{sec:raw-lines}).
\end{itemize}

Change a setting and the pane follows: set the delimiter to semicolon and the marks move to the
semicolons; set \ui{Skip rows} to 2 and the first two lines are struck. It is the quickest way to see why a
file reads wrong, and to see that a change fixes it.

\subsection{The settings}

\begin{center}\small
\begin{tabularx}{\linewidth}{@{}lML@{}}
\toprule
\sffamily\bfseries setting & \sffamily\bfseries what it does & \sffamily\bfseries choices \\
\midrule
\ui{View as} & read the file as a table or as lines of text & auto · table · text \\
\ui{Delimiter} & the column separator & auto · comma · semicolon · tab · pipe · whitespace \\
\ui{Header} & whether the first row names the columns & auto · yes · no \\
\ui{Decimal} & how numbers write their fraction & auto · point (1.5) · comma (1,5) \\
\ui{Encoding} & how the bytes become characters & UTF-8 · Windows-1252 · Latin-1 · UTF-16 LE · UTF-16 BE \\
\ui{Skip rows} & lines to skip before the header & a number \\
\ui{Comment} & skip the leading lines that start with this & a character or word \\
\bottomrule
\end{tabularx}
\end{center}

Changing a setting reads the file again with it, in the same view, at once, and the \ui{Detected} list
follows the new reading. A setting left on \ui{auto} or blank keeps the detected value.

A Datamine file, a Parquet file or a workbook sheet carries its columns and types inside it, so there is
nothing to choose: the popover lists what was found (format and columns; for Parquet, the row groups and
compression), says \ui{a binary table — its columns and types are in the file, so there's nothing here to
choose}, and the settings are greyed out.

\begin{note}
A forced interpretation isn't stored in the index cache: the next time you open the file, lamina detects
it afresh. A lens keeps it (\cref{ch:lens}), and so does the memory of a file's last views
(\cref{sec:last-views}).
\end{note}

\section{When a file can't be shown}\label{sec:failures}

lamina says why rather than showing an empty grid or garbage. The file's tab holds a short note instead
of a table (a \term{note tab}, \cref{sec:note-tabs}):

\begin{description}
\item[\ui{not a table}] the bytes aren't text and no reader recognised them; \menu{View>Raw — the bytes
in hex} shows them (\cref{ch:raw}).
\item[\ui{archive too large}] a Zstandard, xz, bzip2 or tar archive larger than 32~MB: these must be
decompressed whole, in memory (\cref{sec:archives}).
\item[\ui{workbook too large}] a workbook larger than 32~MB, for the same reason (\cref{sec:xlsx}).
\item[\ui{empty archive}], \ui{no data in this workbook} — exactly that.
\item[\ui{couldn't read this …}] a damaged or truncated workbook or archive, with the reader's own
message.
\item[\ui{can't show this Parquet file}] a Parquet file lamina can't decode, with the reason: most
often a compression it doesn't carry (\cref{sec:parquet}).
\end{description}

\begin{tryit}
Watch lamina read a file.
\begin{enumerate}
\item Save these four lines in a text editor as \file{assays.csv}:
\begin{codeblock}
HOLE;FROM;TO;FE
DH1;0,0;1,5;45,2
DH1;1,5;3,0;51,7
DH2;0,0;2,0;38,9
\end{codeblock}
\item Open it in lamina: three rows, columns \col{HOLE}, \col{FROM}, \col{TO}, \col{FE}, the last three
numbers. The semicolon and the decimal comma were both found.
\item Click the kind badge. The pane shows the four lines with the semicolons marked, the first labelled
as the header, and four fields beside each.
\item Set \ui{Delimiter} to \ui{comma}. The pane now marks the decimal commas as delimiters, the field
counts go wrong, and the table behind it changes to match. Set it back to \ui{auto}.
\end{enumerate}
\end{tryit}
